\documentclass{homework}
% \author{Lastname, Firstname}  % Uncomment with your name
\class{CS 2813: Tashfeen's Discrete Structures}
\date{Spring 2023}
\title{Homework 3}
\address{University of Oklahoma, Computer Science}

\begin{document} \maketitle

\question Please read chapters 3 and 4 of Chartrand et al. and write a couple sentences about a topic/example/concept that you found difficult or interesting and why?

\question Consider the following quantified statement: For every real number $x$, there exists a positive real number $y$ such that $y < x^2$.

\begin{enumerate}[label=(\alph*)]
	\item Express this quantified statement in symbols.
	      % \[
		  %     here
	      % \]
	\item Express the negation of this quantified statement in symbols.
	      % \begin{align*}
		  %     here
	      % \end{align*}
	\item Express the negation of this quantified statement in words.
	      % here
\end{enumerate}

\question Prove that if $r$ and $s$ are rational numbers, then $r - s$ is a rational number.

% \begin{proof}
% 	here
% \end{proof}

\question Let $x$ and $y$ be integers. Prove that if $x + y \geq 9$, then either $x \geq 5$ or $y \geq 5$.

% \begin{proof}
% 	here
% \end{proof}

\question Let $m$ and $n$ be two integers. Prove that $mn$ and $m + n$ are both even if and only if $m$ and $n$ are both even.

% \begin{proof}
% 	here
% \end{proof}

\question Disprove: Let $A, B$ and $C$ be sets. If $A \cup B = A \cup C$, then $B = C$.

% here

\question Prove that if $a$ and $b$ are positive real numbers, then $\sqrt{a} + \sqrt{b} \neq \sqrt{a+b}$.

% \begin{proof}
% 	here
% \end{proof}

\question Let $r \geq 2$ be an integer. Prove that $1 + r + r^{2} + \cdots + r^{n} = \frac{r^{n+1}-1}{r-1}$ for every positive integer $n$.

% \begin{proof}
% 	here
% \end{proof}

\question Prove that $\frac{1}{\sqrt{1}} + \frac{1}{\sqrt{2}} + \cdots + \frac{1}{\sqrt{n}} > \sqrt{n+1} $ for every integer $n\geq 3$.

% \begin{proof}
% 	here
% \end{proof}

\question A sequence $a_{1}, a_{2}, a_{3},\cdots$ is defined recursively by $a_{1} = 3$ and $a_{n} = 2a_{n-1} + 1$ for $n\geq 2$.
\begin{enumerate}[label=(\alph*)]
	\item Determine $a_{2}, a_{3}, a_{4},$ and $a_{5} $.
	\item Based on the values obtained in (a), make a guess for a formula for $a_{n}$ for every positive integer $n$ and use induction to verify that your guess is correct.
	      % \begin{proof}
		  %   here
	      % \end{proof}
\end{enumerate}

\question In Example 4.36, we saw that $n^{th}$ Fibonacci number $F_{n} \leq 2^{n}$. Prove that $F_{n}\leq (\frac{5}{3})^{n}$ for every positive integer $n$.

% \begin{proof}
% 	here
% \end{proof}

\question A sequence \{$a_{n}$\} is defined recursively by $a_{1} = 5$, $a_{2} = 7$ and ${a_{n} = 3a_{n-1}-2a_{n-2}-2}$ for $n\geq3$. Prove that $a_{n} = 2n + 3$ for every positive integer $n$.

% \begin{proof}
%   here
% \end{proof}

\end{document}
